\documentclass[10pt,a4paper]{amsart}
\usepackage[margin=19mm]{geometry}
\usepackage{amsmath,amssymb,mathtools}
\usepackage[scr=rsfs,cal=euler]{mathalfa}
\usepackage{xfrac}
\usepackage[unicode,hidelinks,bookmarks=false]{hyperref}
\newcommand{\ie}{i.e.\ }
\newcommand{\cf}{cf.\ }
\newcommand{\resp}{respectively\ }
\newcommand{\Subsection}{\subsection}
\providecommand{\nobreakdash}{\nobreak\hbox{-}}
\setlength{\emergencystretch}{2em}
\multlinegap0pt
\usepackage{amssymb}
\newtheorem{lemma}{Lemma}
\newtheorem{prop}{Proposition}
\newcommand{\E}{\mathbb{E}}
\newcommand{\Z}{\mathbb{Z}}
\title{EFFECTIVE RADIUS OF A DISCRETE MOVING POLYMER}
\author{Jiaming Chen}
\date{Source: arXiv:2507.22248v2 (CC BY 4.0)}
\begin{document}
\maketitle

\noindent\emph{Source and licence.} EFFECTIVE RADIUS OF A DISCRETE MOVING POLYMER. Version arXiv:2507.22248v2 (CC BY 4.0), distributed by arXiv under the Creative Commons Attribution 4.0 International licence, http://creativecommons.org/licenses/by/4.0/, as recorded in the arXiv licence metadata for this exact version and shown on the abstract page https://arxiv.org/abs/2507.22248v2. Full licence notice: \texttt{licenses/arxiv-2507.22248-CC-BY-4.0.txt}.\par\smallskip

\noindent\textit{Editorial scope.} Counted unit: the article's prop lower (source0.tex lines 543--549). The article's complete original proof of it is delivered with it (source0.tex lines 550--618, one \texttt{proof} environment of 69 lines). No mathematics is written, repaired or re-derived: the statement and the proof are the article's own lines, in the article's own notation, and no retained line is substituted or re-flowed.  Retained uncounted as the source-local context the proof needs: lines 135--140; lines 153--156; lines 179--182; lines 184--187; lines 246--250.  Nothing was omitted from any retained span, so the package carries no omission record.  No retained line contains a citation, so the wrapper carries no bibliography and no bibliography entry.  No retained line contains a figure environment, an image-inclusion command or drawing code, so no image is delivered and the package declares no asset dependency.  Editorial formatting only, in addition: an AMS \texttt{amsart} preamble that restates the article's own declarations and macros used by the retained lines, namely the environments \texttt{lemma}, \texttt{prop} on separate counters, together with the standard packages those lines need.  The retained environments print in this document as Lemma 1, Proposition 1 in order of appearance, whereas the article prints the same items with the numbering of its own full document.  The complete native source of the article as distributed in the arXiv e-print for this exact version is bundled unchanged as \texttt{sources/182401/source0.tex}.
\begin{equation}
\label{localTime}
\begin{split}
    N_d(t):=\sum_{0\leq i\neq j\leq J-1}1_{\vert u(t,i)-u(t,j)\vert\leq d}.
    \end{split}
\end{equation}

\begin{equation}
\label{Qmeas}
Q_{T,J,\beta,d,\alpha}(A) =\frac{ \mathbb{E}^{P_{T,J}}[1_{A}\mathcal{E}_{T,J,\beta,d,\alpha}]}{Z_{T,J,\beta,d,\alpha}}, \text{ for all measurable set $A$},
\end{equation}

\begin{equation}
\label{ubar}
\bar{u}(t) = \frac{1}{J}\sum_{n=0}^{J-1}u(t,n),
\end{equation}

\begin{equation}
\label{RTJ}
R(T,J) = \left[\frac{1}{TJ}\sum_{t=1}^T\sum_{n=0}^{J-1}(u(t,n)-\bar{u}(t))^2\right]^{1/2}.
\end{equation}

\begin{lemma}\label{ZT} There exist positive constants $C,C',\tilde{d}$ and $\tilde{J}$ such that, for any $J>\tilde{J}$ and $\frac{\tilde{d}J}{\log J}\geq d\geq 1$, we have 
    \[\frac{1}{T}\log Z_T\geq \begin{cases}
            -C\beta^{\frac{2}{3}}d^{\frac{4-2\alpha}{3}}J^{\frac{1}{3}}(\log J)^{\frac{2}{3}} & \beta\geq C'd^{1+\alpha}J^{-2}(\log J)^{-1},\\
            -Cd^2J^{-1}&\beta<C'd^{1+\alpha}J^{-2}(\log J)^{-1}.\end{cases}\]
\end{lemma}

\begin{prop}
\label{lower}
    There exist positive constants $\tilde{d}, \tilde{J}, \beta_C$ and $K_1$, such that for $J>\tilde{J}$, $\beta \geq\beta_Cd^{4+\alpha}J^{-2}(\log J)^{2}$ and $\frac{\tilde{d}J}{\log J}\geq d\geq 1$,
    \begin{equation*}
    \lim_{T\to \infty}Q_T\left[R(T,J)\leq K_1\beta_C^{\frac{1}{3}}d J\right]=0.
\end{equation*}
\end{prop}

\begin{proof}
Fix $K>0$ and define the event $A_{K,T}:=\{R(T,J)\leq K\}$. By the definition of $R(T,J)$ in \eqref{RTJ}, the event $A_{K,T}$ implies
\begin{equation}
\label{inclusion1}
    \left\vert \left\lbrace t\in \mathcal{T}:\frac{1}{J}\sum_{n=0}^{J-1}(u(t,n)-\bar{u}(t))^2\leq2K^2\right\rbrace\right\vert\geq \frac{T}{2},
\end{equation}
where $|\cdot|$ denotes cardinality and $\bar u(t)$, defined in \eqref{ubar}, is the spatial average of $u(t,\cdot)$ at time $t$. Next observe that if
\begin{equation*}
    \frac{1}{J}\sum_{n=0}^{J-1}(u(t,n)-\bar{u}(t))^2\leq2K^2,
\end{equation*}
then a simple counting argument shows that
\begin{equation}
    \label{inclusion2}
    \left\vert \left\lbrace n\in \mathcal{J}:u(t,n)\in[\bar{u}(t)-2K,\bar{u}(t)+2K]\right\rbrace\right\vert\geq \frac{J}{2}.
\end{equation}
Set $d_t^{\pm}=\bar{u}(t)\pm 2K$, so that $d_t^+-d_t^-=4K$. There exist integers $z^+,z^-\in\mathbb{Z}$ and a parameter $\gamma\in[0,1]$ such that
\begin{equation*}
    d_t^+=z^+d-\gamma d, \quad d_t^-\in(z^-d-\gamma d,z^-d+(1-\gamma)d].
\end{equation*}
To measure the local density of the profile $u(t,\cdot)$ at time $t$, we define the number of spatial sites that fall within a given interval of length $d$ as
\begin{equation}
\label{localT}
    l_{d,t}(z)=\sum_{i=0}^{J-1}1_{u(t,i)\in(zd-\gamma d,zd+(1-\gamma)d]}.
\end{equation}

Combining \eqref{inclusion1} and \eqref{inclusion2}, we see that on the event $A_{K,T}$,
\begin{equation}
\label{localbound}
    \left\vert \left\lbrace t\in \mathcal{T}:\sum_{z=z^-}^{z^+-1}l_{d,t}(z)\geq \frac{J}{2}\right\rbrace\right\vert\geq \frac{T}{2}.
\end{equation}
Note that even if two values $u(t,i)$ and $u(t,j)$ satisfy $\vert u(t,i)-u(t,j)\vert<d$, they need not belong to the same interval in the partition \eqref{localT}. Nevertheless, the following lower bound remains valid:
\begin{equation}
\label{localTineq}
    N_d(t)\geq\sum_{z\in\Z}l^2_{d,t}(z)-J,
\end{equation}
where $N_d(t)$ is the approximate self-intersection count defined in \eqref{localTime}.

Applying the Cauchy-Schwarz inequality to the sum over $z$, and then using \eqref{localbound} and \eqref{localTineq}, yields on $A_{K,T}$
\begin{equation}
\begin{split}
\label{QmeasB}
\sum_{t=1}^{T}N_d(t)&\geq\sum_{t=1}^{T}\sum_{z\in\Z}l^2_{d,t}(z)-TJ\geq\sum_{t=1}^{T}\sum_{z=z^-}^{z^+-1}l^2_{d,t}(z)-TJ\\
    &\geq \frac{1}{z^+-z^-}\sum_{t=1}^{T}\left(\sum_{z=z^-}^{z^+-1}l_{d,t}(z)\right)^2-TJ\\
    &\geq \frac{d}{4K+d}\cdot\frac{T}{2}\cdot\left(\frac{J}{2}\right)^2-TJ = TJ\left(\frac{d J}{32K+8d}-1\right).
    \end{split}
\end{equation}
At this point it is useful to note that $K$ can be at most a constant multiple of $dJ$. Indeed, otherwise one could take $u(t,i)=id$, in which case the radius of gyration is of order $dJ<K$, while $N_d(t)=0$ for every $t$. In the continuous model, the diagonal set has zero measure, whereas in the discrete model the diagonal consists of exactly $J$ points and therefore contributes nontrivially to the counting measure. To obtain a lower bound for $K$ of at least order $J$, one must therefore take $d>1$ sufficiently large so that the diagonal contribution does not overwhelm the off-diagonal self-interaction terms.

We now estimate the logarithmic asymptotic probability under $Q_T$. Using Lemma \ref{ZT}, together with \eqref{QmeasB} and the definition of the penalizing factor $\mathcal{E}_T$ in \eqref{Qmeas}, we find
\begin{equation}
\label{smalldistance}
\begin{split}
    &\limsup_{T\to \infty}\frac{1}{T}\log Q_T\left[R(T,J)\leq K\right]\\
    &\leq\limsup_{T\to \infty}\frac{1}{T}\log \E^{P_T}\left[\mathcal{E}_T1_{R(T,J)\leq K}\right]-\liminf_{T\to\infty}\frac{1}{T}\log Z_T\\
    &\leq\begin{cases}
            -\frac{\beta J}{d^{\alpha}}\left(\frac{d J}{32K+8d}-1\right)+C\beta^{\frac{2}{3}}d^{\frac{4-2\alpha}{3}}J^{\frac{1}{3}}(\log J)^{\frac{2}{3}} & \beta\geq C'd^{1+\alpha}J^{-2}(\log J)^{-1},\\
            -\frac{\beta J}{d^{\alpha}}\left(\frac{d J}{32K+8d}-1\right)+Cd^2J^{-1}&\beta<C'd^{1+\alpha}J^{-2}(\log J)^{-1}.\end{cases}
    \end{split}
\end{equation}
If $\beta<C'd^{1+\alpha}J^{-2}(\log J)^{-1}$, then $\frac{\beta J}{d^\alpha}<C'd J^{-1}(\log J)^{-1}<Cd^2J^{-1}$, so for large $J$ the upper bound of \eqref{smalldistance} cannot be negative. By contrast, if $\beta \geq\beta_Cd^{4+\alpha}J^{-2}(\log J)^{2}$ and we choose $K=K_1\beta_C^{\frac{1}{3}}d J$, then \eqref{smalldistance} becomes
\begin{equation*}
    \limsup_{T\to \infty}\frac{1}{T}\log Q_T\left[R(T,J)\leq K_1\beta_C^{\frac{1}{3}}d J\right]\leq \left(-\frac{\beta_C^{1/3}d J}{32K_1\beta_C^{\frac{1}{3}}d J+8d}+C'\right)\beta^{\frac{2}{3}}d^{\frac{4-2\alpha}{3}}J^{\frac{1}{3}}(\log J)^{\frac{2}{3}}.
\end{equation*}
Choosing $K_1$ sufficiently small makes the coefficient in parentheses strictly negative. Hence the exponential decay rate is negative, and therefore
\begin{equation*}
    \lim_{T\to \infty}Q_T\left[R(T,J)\leq K_1\beta_C^{\frac{1}{3}}d J\right]=0,
\end{equation*}
which completes the proof.
\end{proof}

\end{document}
